\section{Response to Reviewer \texttt{SDFf}}

\subsection*{Summary}

The paper proposes a three-stage pipeline that (i) learns low-variance outcome models on pooled source RCT IPD (proxy labels), (ii) debiases the baseline using scarce target placebo (gold) outcomes via a sparse linear correction (LASSO), and (iii) embeds the anchored outcome models into a doubly robust (DR) / orthogonal learner to estimate patient-level CATEs. However, several of the core assumptions and statistical-rate claims appear underspecified or overly optimistic for the data regime under consideration (e.g., enough gold samples for cross-validation and estimation, bounded covariate shift without guaranteed bounded outcome regression, and strong assumptions on the transportability of a sparse covariate subset across sites). The proposed method may be promising in practice, but the theoretical conditions currently feel stronger than what the empirical setting can support.

The manuscript would be strengthened by additional experiments (e.g., non-linear treated-arm corrections, non-linear data-generating processes, comparisons against existing transport or meta-analytic methods), deeper evaluation of experiments via many Monte Carlo realizations (instead of a single synthetic data generation realization), and statistical hypothesis testing to more rigorously demonstrate performance gains in the proposed approach.


\subsection*{Summary of Strengths}

\begin{enumerate}[label=\arabic*.]

\item The paper clearly articulates a practical and understudied problem of estimating treatment effects in a target population with only placebo outcomes available, using multi-trial RCT IPD as proxy data. The ``proxy--gold'' formalism provides a nice decomposition of the treatment effect into two terms: one term from the proxy outcomes (external trials) and one term from the gold outcomes (internal targeted trial).

\item The ablation studies clearly show the contribution of each stage: proxy-only, anchor-only, and proxy+anchor (full model), which demonstrate where improvements are driven from for CATE estimation.

\item The modular three-stage design enables flexible learners, where Stage~1 permits any high-capacity learner (e.g., random forest, neural network), while Stages~2 and~3 attempt to recover causal validity via calibration to target placebo and anchored treated outcomes.

\item The manuscript lays out comprehensive empirical robustness check concepts that illustrate practical failure modes and build intuition about where the method succeeds or fails.

\end{enumerate}

\subsection*{Summary of Weaknesses}

\subsubsection{Conditional exchangeability and unmeasured effect modifiers.}  
A core requirement in transport theory is conditional exchangeability $Y(0), Y(1) \perp C|X$ i.e., that outcomes are independent of the trial indicator $C$ given covariates $X$. The paper explicitly states that this does not hold, as outcome models from source trials may be biased for the target and covariate distributions drift across sites. This raises concerns that covariate shift may be correlated with unmeasured effect modifiers (e.g., sites recruiting more severe patients may also differ in unmeasured clinical protocols). It is unclear mathematically whether the sparse bias regression term is sufficient to address this issue.\\

{\color{blue}We thank the reviewer for highlighting this fundamental issue and fully agree with the concern. We emphasize at the outset that our method does not rely on, nor does it claim to recover, classical conditional exchangeability across sites. In particular, we do not assume 
\[
(Y(0),Y(1)\perp C | X)
\]
nor do we claim nonparametric identification of target treatment effects in the presence of unmeasured effect modifiers.

To clarify the scope of our contribution, it is useful to distinguish baseline risk transportability from treatment-effect transportability. Existing multi-site transport and meta-analytic methods typically require assumptions such as:
\begin{itemize} 
\item Exchangeability over sites \citep{dahabreh2023efficient}: $$Y(a)\perp C | X,$$
\item Sample ignorability for treatment effects 
\citep{hong2025estimating}: $$C \perp (Y(1)-Y(0)) | X.$$ 
\end{itemize}
Both of which rule out unmeasured effect modifiers. As noted by \cite{rott2025causally}, these assumptions are often restrictive and implausible in practice.

Our approach explicitly departs from this identification paradigm. Rather than assuming transport bias is zero, we treat it as present but structured. Specifically, we leverage observed target-site placebo outcomes to anchor and recalibrate baseline risk. This anchoring step operates entirely through $Y(0)$ and therefore addresses baseline risk miscalibration, not unobserved treatment-effect heterogeneity.

Importantly, we do not claim that placebo anchoring recovers bias arising from unmeasured effect modifiers that differentially affect treated and untreated outcomes. If such modifiers induce treatment-specific deviations beyond baseline risk, they remain fundamentally unidentifiable without treated outcomes in the target population. Our estimator should therefore be interpreted as a calibrated, model-based transport estimator, not a design-based identification result.

Assumption A5 formalizes this distinction. It is introduced as a regularity condition on the complexity of cross-site deviations in the outcome regression, not as a substitute for conditional exchangeability. When A5 holds approximately, anchoring improves calibration and estimation stability; when it is violated, performance degrades smoothly toward proxy-only baselines rather than exhibiting a sharp identification failure (see response to Q5 further below regarding concerns about A5).

The manuscript has been revised to make this point explicit and to avoid language suggesting that unmeasured effect modification is resolved. We emphasize that our contribution lies in robust calibration under acknowledged transport bias, rather than causal identification under strong exchangeability assumptions.}


\subsubsection{Total variation bound and outcome stability.}  The argument that $\mu_a(x)$ is the same between sites because the total variation (TV) distance of covariates distributions across trials (sites) is bounded is questionable. Also, the total variation between two probability density functions is always less than or equal to $1$, so why is it bounded by infinity in A4? Translating TV distance bounds into bounds on the predicted outcome of $\mu_a$ requires further smoothness or Lipschitz assumptions on $\mu_a(x)$, which are not discussed. For example (with a different distribution measure), if $\mu_a$ is Lipschitz with constant $L$, then
\[
W_1(\mu_{\sharp} P, \mu_{\sharp} Q) \le L W_1(P, Q),
\]
where $W_1$ is the $1$-Wasserstein metric and $\mu_{\sharp}$ is the pushforward operation on the input covariate distribution \citep{villani2008optimal}.


{\color{blue}
We thank the reviewer for carefully pointing out the issues with our original formulation of Assumption~A4. We agree that the original statement was imprecise and potentially misleading, and we apologize for the sloppiness in its presentation.

In the original manuscript, Assumption~A4 was stated as a bound on covariate shift across sites using total variation distance:
\[
\textbf{(Original A4)} \qquad 
d_{\mathrm{TV}}\!\left(P_{X \mid c_1}, P_{X \mid c_2}\right) \le \Delta < \infty,
\]
for all sites $c_1, c_2$, where $d_{\mathrm{TV}}$ denotes total variation distance. 

As the reviewer correctly notes, total variation distance is always bounded above by~1, and therefore the condition $\Delta < \infty$ is redundant. 
More importantly, as pointed out by the reviewer, bounding covariate shift in total variation does not, by itself, imply stability of outcome regressions or bounded prediction error without additional smoothness assumptions. 
We further clarify that we never intended to assume or claim that $\mu_a(x)$ is identical across sites; rather, the original assumption was meant to impose weak regularity for estimation, not equality or transportability.


To address this concern, we have revised Assumption~A4 to a Lipschitz-type regularity condition on the conditional outcome regressions, removing total variation distance entirely. The revised assumption is stated as follows:
\[
\textbf{(Revised A4)} \qquad
\left| \mu_{a,c}(x) - \mu_{a,c}(x') \right|
\;\le\; L \,\|x - x'\| ,
\quad \forall a \in \{0,1\},\ \forall c,\ \forall x,x',
\]
where $\mu_{a,c}(x) = \mathbb{E}[Y \mid A=a, X=x, c]$ and $L < \infty$ is a uniform Lipschitz constant.

This revised assumption serves solely as a \emph{regularity condition} ensuring that moderate covariate shift across sites does not induce arbitrarily large changes in predicted outcomes for nearby covariate profiles.
It does \emph{not} imply exchangeability, transportability, or stability of treatment effects across sites, and it is not used to justify causal identification.

All structural assumptions regarding how cross-site differences manifest in outcome regressions are now explicitly delegated to Assumption~A5, which we frame as a regularity condition enabling calibrated estimation rather than as an identification assumption. We believe this revision clarifies the role of Assumption~A4 and directly addresses the reviewer’s concern.
}

\subsubsection{Monte Carlo evaluation and statistical testing.}  
For the synthetic experiments, there should be many Monte Carlo realizations to compute means and standard deviations of PEHE, MAE of ATE, calibration RMSE for $\mu_0$ and $\mu_1$ in both ablation and robustness checks. Otherwise, it is unclear whether observed improvements are due to randomness in data generation. Statistical hypothesis testing (e.g., Friedman test, with causal variants as treatments and metrics as blocks) is suggested. The causal variants in table II are the ``treatments" and the metrics are the ``blocks". Lastly you could vary the number of covariates used.\\

{\color{blue} We agree with the reviewer and have substantially expanded the Monte Carlo evaluation. All synthetic experiments are now run with 100 independent replications, and we report mean $\pm$ standard deviation for all metrics, including PEHE, absolute ATE error, Spearman correlation, policy regret, and CATE calibration error (ECE). The revised experiments section includes the following comprehensive evaluations:
\begin{itemize}
\item \textbf{Target budget vs.\ dimensionality sweep:} We vary target sample sizes $(m_0, m_1)$ across five budget regimes and covariate dimension $p \in \{10, 20, 50, 100\}$, evaluating all 10 comparison methods.
\item \textbf{Source site scaling:} We vary the number of source sites $C \in \{2, 5, 10, 20, 50\}$ while holding the target budget fixed.
\item \textbf{A5 violation sensitivity:} We sweep sparsity ratio and nonlinearity $\Gamma$ to quantify graceful degradation when Assumption A5 is violated.
\end{itemize}
Statistical comparison across methods is facilitated by the comprehensive error bars and an average rank summary table that aggregates performance across all conditions and metrics (see revised manuscript Table~III and Appendix~C).}

\subsubsection{Claim of estimation without a treated arm in the target.}  
There is a claim that treatment effects are estimated in the target population with no treated arm (first page in the introduction), but if so, how is $\beta_{1,0}$ estimated? One may transport support but not estimated coefficients. Later in the paper, it is stated that both arms are present in the gold data.

{\color{blue}
We thank the reviewer for pointing out this ambiguity. We agree that our original wording could be misread as claiming that we can \emph{estimate} the treated-arm regression in the target from target placebo data alone, or that we can recover a target treated coefficient vector $\beta_{1,0}$ without treated outcomes. This is not what our updated (implemented) procedure does. We have revised the introduction and estimator description to make the distinction explicit.

\paragraph{No identification of the target treated arm from placebo-only target data.}
When the target contains placebo outcomes only ($A\equiv 0$ in $S=0$), the treated regression $\mu_{1,0}(x)$ and hence the true target CATE $\tau_0(x)=\mu_{1,0}(x)-\mu_{0,0}(x)$ are not nonparametrically identified from target data alone. We make no claim to the contrary, and we no longer describe the disconnected-target setting as ``estimating'' $\beta_{1,0}$.

\paragraph{What we do instead in disconnected targets (Option B / \texttt{Proposed-B}).}
In disconnected targets, we do \emph{not} attempt to estimate a target treated coefficient vector $\beta_{1,0}$. Our goal is prediction/transport of a CATE function using source trials where both arms are observed, with the target placebo arm used only to select credible sources. Concretely, the implemented \texttt{Proposed-B} procedure is:
\begin{enumerate}[leftmargin=*]
\item \textbf{Placebo-based screening.} Using only the target placebo sample $\mathcal{D}_{0,0}=\{(X,Y):S=0,A=0\}$, we run placebo-arm transfer screening (Trans-GLM detection) to select a subset of sources $\widehat{\mathcal{C}}_0$ whose placebo outcome models are compatible with the target on the target covariate support.
\item \textbf{Source DR learning.} On the selected sources $\bigcup_{c\in\widehat{\mathcal{C}}_0}\mathcal{D}_c$ (both arms observed), we fit outcome regressions and construct doubly robust pseudo-outcomes, then learn a CATE predictor $\widehat{\tau}_B(x)$ \emph{on the selected sources}.
\item \textbf{Transport by prediction.} We evaluate $\widehat{\tau}_B(x)$ on target covariates $X\sim P_0$ and output $\widehat{\tau}_B(x)$ as the transported CATE estimate for target individuals.
\end{enumerate}
Thus, in Option B we transport a \emph{function} $\tau(\cdot)$ learned from screened sources; we do not claim recovery of target treated coefficients or a target treated regression from placebo-only target data.

\paragraph{Role of Assumption~A6 (screening-valid transportability).}
The relevant working-model condition for Option B is Assumption~A6, which states that placebo-based screening identifies a subset of sources whose CATE functions are close to the target CATE on the target covariate support (up to an explicit approximation level $\epsilon_\tau$). Under A6, our guarantee is a \emph{transport error decomposition} (estimation on selected sources $+$ structural transport term $\epsilon_\tau$ $+$ screening error), not identification of $\beta_{1,0}$ or $\mu_{1,0}$ from target placebo alone.

\paragraph{Clarification of ``gold data'' and presence of both arms.}
We have revised the manuscript to avoid conflating (i) the placebo-only target regime (Option B), where the target gold data consist only of placebo outcomes, with (ii) the connected-target regime (Option A / \texttt{Proposed-CF}), where both arms are observed in the target and the DR learner is fit on target data. References to ``gold'' now explicitly indicate whether they refer to target placebo only (for anchoring/screening) or to a target dataset containing both arms (used only in Option A and in benchmarking experiments).
}

\subsubsection{Restrictive linear and sparse bias model.}  
The model constrains the bias term to be linear and sparse, which is restrictive and not well justified theoretically. If the true correction term were nonlinear or non-sparse, it is unclear how the method would perform. Sensitivity analyses are suggested for violations of Assumption~A5 (non-sparse or nonlinear bias).

{\color{blue}
We apologize for the lack of clarity in the original presentation and thank the reviewer for highlighting an important practical concern. In particular, we regret that Assumption~A5 was stated to be an identification assumption. 
We would like to clarify that Assumption~A5 is \emph{not} intended to guarantee full causal identification or transportability of treatment effects. 
Rather, it is a \emph{structural regularity condition} that restricts the class of admissible cross-site deviations and enables stable estimation and calibrated extrapolation when treated outcomes are unavailable in the target population.

Assumption~A5 is motivated by a growing body of transfer learning and clinical prediction literature showing that differences between study populations are often driven by a limited number of clinically meaningful factors, rather than by wholesale changes in the outcome-generating mechanism. In practice, site- or population-level differences in randomized trials frequently arise from shifts in baseline risk calibration, prevalence of key biomarkers, disease severity distributions, or measurement and enrollment practices, rather than from entirely different treatment-response mechanisms.

In the proxy--gold transfer sparse regression based learning literature, \cite{bastani2021predicting} formalizes this idea by modeling discrepancies between abundant proxy data and scarce gold-standard data as low-dimensional corrections to an otherwise well-learned prediction model. Importantly, this assumption is supported empirically in healthcare prediction tasks, where sparse adjustments are often sufficient to reconcile differences between external datasets and target clinical populations.

Crucially, Assumption~A5 does \emph{not} require that transport bias be exactly sparse or perfectly linear.
As discussed in \cite{bastani2021predicting}, the bias term recovered in the second stage of a two-step transfer estimator typically consists of a dominant low-dimensional signal combined with residual noise propagated from the proxy model. Established high-dimensional theory shows that LASSO-based procedures remain statistically valid under \emph{approximate sparsity}, with estimation error decomposing into a variance component from the proxy stage and a regularization error term. \cite{buhlmann2011statistics, belloni2012sparse, bastani2021predicting} further show that this structure yields meaningful sample-efficiency gains while degrading gracefully when sparsity is only approximate.

Importantly, approximate sparsity does not restore identification of the true target CATE.
When Assumption~A5 is violated, the estimator converges to a different working-model target and incurs a non-vanishing approximation bias, which we explicitly characterize in our theory and robustness experiments.

From a clinical standpoint, this assumption aligns closely with standard practices in individual participant data meta-analysis and model transportability. These literatures emphasize parsimonious modeling of effect modification across trials, often focusing on a small number of known or suspected modifiers, rather than unrestricted interaction structures \citep{riley2021using,dahabreh2023efficient}.
Similarly, the clinical prediction model updating literature consistently finds that cross-site differences are most effectively addressed through recalibration or limited coefficient revision, rather than full refitting of complex models \citep{wynants2017key, van2019calibration, steyerberg2019validation}.

This sparsity-of-contrast perspective has also been formalized in modern multi-source transfer learning theory. \cite{li2022transfer} study settings with multiple auxiliary datasets and explicitly assume that informative sources differ from the target through sparse contrast vectors, while allowing non-informative or heterogeneous sources to coexist. 
They establish minimax-optimal rates and robustness guarantees under this assumption, directly reflecting the realities of multi-center clinical trial data.
\cite{tian2023transfer} extend these results to high-dimensional generalized linear models, showing that sparse contrasts are sufficient for improved prediction and estimation, while violations can lead to negative transfer—an observation that mirrors clinical experience when inappropriate external data are over-weighted.

Finally, we emphasize that Assumption~A5 does \emph{not} assert that covariate distributions themselves shift sparsely across sites. That is, we do not assume for sites $c,c'$
\[
\|X_{c} - X_{c'}\|_0 \le q \quad \text{for small } q \in \mathbb{N}.
\]
Rather, the assumption concerns the \emph{outcome regression}, namely that differences in how patient characteristics map to outcomes across sites can be well approximated by sparse contrasts in regression coefficients. Under a linear approximation,
\[
Y_{c_i} = \beta_{c}^\top X_{c} + \varepsilon_{c}, \qquad
Y_{c_j} = \beta_{c'}^\top X_{c'} + \varepsilon_{c'},
\]
we assume
\[
\|\beta_{c} - \beta_{c'}\|_0 \le q \quad \text{for small } q \in \mathbb{N}.
\]
This reflects the clinical belief that only a small subset of patient features—such as key risk factors or effect modifiers—drive systematic differences in outcomes across trial populations.

\textbf{Accordingly, Assumption~A5 is now stated as a regularity condition for estimation rather than an identification assumption.} When the assumption holds approximately, placebo anchoring provides informative calibration signal and improves estimation efficiency; when it is violated, performance degrades gracefully toward proxy-only baselines rather than exhibiting a sharp identification failure.

}


\subsubsection{Shared support across treatment arms.}  
It is unclear why it is necessarily true that different treatment arms share the same support set of covariates.

{\color{blue}
We thank the reviewer for raising this point, as it highlights an important distinction between biological assumptions and design-based assumptions. Shared support across treatment arms is not a biological claim about treatment response, but a standard \emph{design property} of randomized controlled trials.

In our setting, treatment assignment occurs after enrollment and is randomized with known probabilities. Consequently, for any baseline covariate pattern observed at enrollment, each participant has a positive probability of being assigned to either treatment arm. This corresponds to the standard positivity (or overlap) condition in causal inference for randomized experiments, under which baseline covariates have overlapping support across treatment groups by design \citep{imbens2015causal, hernan2010causal, dahabreh2023efficient}

We emphasize that this assumption applies only to \emph{baseline covariates measured prior to treatment assignment}. It does not require that post-treatment variables, intermediate outcomes, or biologically driven response mechanisms share support across arms. Nor does it assert that treatment effects are homogeneous across covariate values. Rather, it ensures that conditional expectations such as $\mathbb{E}[Y \mid A=a, X=x]$ are well-defined for both arms over the observed covariate space.

We also acknowledge that violations of this assumption can occur in practice, for example due to deterministic treatment assignment, protocol-driven exclusions after enrollment, or conditioning on post-baseline variables. Such violations would indeed invalidate the overlap condition and, consequently, the proposed approach. We now explicitly clarify this distinction, scope, and failure mode in the revised manuscript.
}

\subsubsection{Disconnected trial networks.}  
The paper motivates disconnected trial networks, but does not explicitly evaluate performance on disconnected networks (e.g., no A--B or B--C comparisons, only A--C gaps). For the synthetic dataset, it is unclear how disconnected networks are constructed if estimation requires treated arms. Clarification is needed.\\

{\color{blue} 

We thank the reviewer for noting that the original manuscript motivated \emph{disconnected trial networks} but did not explicitly evaluate performance in a setting where the target site has no treated outcomes. We agree with this concern and have revised the manuscript to (i) clarify the data structure at each site, (ii) explain how estimation proceeds without treated target data, and (iii) include an explicit experimental evaluation.

\medskip
\noindent
\textbf{Data structure in the disconnected regime.}
We define disconnectedness operationally: the target site observes \emph{placebo outcomes only}, with no treated samples available for estimation. The data structure is as follows:

\begin{center}
\begin{tabular}{lcc}
\hline
\textbf{Site} & \textbf{Placebo arm} & \textbf{Treated arm} \\
\hline
Source 1 & \checkmark & \checkmark \\
Source 2 & \checkmark & \checkmark \\
$\vdots$ & $\vdots$ & $\vdots$ \\
Source $C$ & \checkmark & \checkmark \\
\textbf{Target} & \checkmark & $\times$ \\
\hline
\end{tabular}
\end{center}

In this setting, there is no direct observation of $Y(1)$ in the target population. Standard doubly robust learners (e.g., TargetOnlyDR) and methods requiring target treated outcomes (e.g., AnchorOnly, Proposed Option A) are \emph{undefined}. This is precisely the motivating regime for our framework.

\medskip
\noindent
\textbf{How estimation proceeds without treated target data (Option B).}
When $m_1 = 0$, our estimator operates as follows:
\begin{enumerate}
\item \textbf{Stage 1 (Proxy learning):} Learn outcome regressions $\widehat{\mu}_0^{\text{proxy}}, \widehat{\mu}_1^{\text{proxy}}$ from pooled source RCT data where both arms are observed.
\item \textbf{Stage 2 (Placebo anchoring):} Using target placebo outcomes $(X_i, Y_i)$ with $A_i=0$, estimate a sparse correction $\widehat{\beta}_{0,0}$ by regressing residuals $Y_i - \widehat{\mu}_0^{\text{proxy}}(X_i)$ on covariates via LASSO.
\item \textbf{Cross-arm transfer (Assumption A6):} From source sites, learn a low-rank operator $\widehat{M}$ relating placebo and treated corrections. Construct the target treated correction as $\widehat{\beta}_{1,0} = \widehat{M} \widehat{\beta}_{0,0}$.
\item \textbf{Stage 3 (CATE estimation):} Form anchored outcome models $\widehat{\mu}_{a,0}^{\text{anch}}(x) = \widehat{\mu}_a^{\text{proxy}}(x) + x^\top \widehat{\beta}_{a,0}$ and construct DR pseudo-outcomes using \emph{source data} (where both arms exist) to estimate the target CATE.
\end{enumerate}

This procedure yields a \emph{working-model} CATE estimate for the target population. It does not claim nonparametric identification from target data alone, but provides calibrated extrapolation under the structural assumptions A5--A6.

\medskip
\noindent
\textbf{Explicit disconnected-regime experiment.}
We evaluate all methods that remain well-defined when $m_1 = 0$: ProxyOnly, AnchorPlugin, IPWTransport, OM-Transport, EntropyBalancing, and Proposed-B. Methods requiring treated target outcomes (TargetOnlyDR, AnchorOnly, Proposed, Proposed-CF) are omitted as they are undefined in this regime.

Table~\ref{tab:disconnected-m10} presents results for the disconnected regime across covariate dimensions $p \in \{10, 20, 50, 100\}$ with $m_0 = 50$ target placebo observations and $m_1 = 0$ (no target treated outcomes). Methods requiring treated target outcomes---TargetOnlyDR, AnchorOnly, Proposed (Option A), and Proposed-CF---are undefined in this regime and therefore omitted.

\begin{table}[t]
\centering
\tiny
\setlength{\tabcolsep}{2pt}
\renewcommand{\arraystretch}{1.1}
\begin{tabular}{@{}l|cccc|cccc|cccc@{}}
\hline
& \multicolumn{4}{c|}{\textbf{PEHE$\downarrow$}} & \multicolumn{4}{c|}{\textbf{ATE$\downarrow$}} & \multicolumn{4}{c}{\textbf{ECE$\downarrow$}} \\
\textbf{Method} & 10 & 20 & 50 & 100 & 10 & 20 & 50 & 100 & 10 & 20 & 50 & 100 \\
\hline
ProxyOnly & 2.3{\tiny$\pm$0.7} & 3.4{\tiny$\pm$0.8} & 5.8{\tiny$\pm$1.2} & 8.2{\tiny$\pm$1.2} & 0.7{\tiny$\pm$0.5} & 0.9{\tiny$\pm$0.7} & 1.3{\tiny$\pm$1.0} & 1.9{\tiny$\pm$1.6} & 0.9{\tiny$\pm$0.5} & 1.1{\tiny$\pm$0.6} & 1.6{\tiny$\pm$0.9} & 2.2{\tiny$\pm$1.4} \\
AnchorPlugin & 1.7{\tiny$\pm$0.7} & 2.4{\tiny$\pm$0.9} & 4.5{\tiny$\pm$1.4} & 6.9{\tiny$\pm$1.2} & \textbf{0.5{\tiny$\pm$0.5}} & \textbf{0.6{\tiny$\pm$0.5}} & 0.9{\tiny$\pm$0.7} & 1.3{\tiny$\pm$1.0} & \textbf{0.6{\tiny$\pm$0.4}} & \textbf{0.7{\tiny$\pm$0.4}} & 1.1{\tiny$\pm$0.6} & 1.7{\tiny$\pm$0.8} \\
IPWTransport & 1.8{\tiny$\pm$0.8} & 2.2{\tiny$\pm$1.1} & 3.3{\tiny$\pm$1.8} & 4.1{\tiny$\pm$2.0} & 0.8{\tiny$\pm$0.6} & 0.8{\tiny$\pm$0.6} & 0.8{\tiny$\pm$0.6} & 1.0{\tiny$\pm$0.9} & 0.8{\tiny$\pm$0.5} & 0.9{\tiny$\pm$0.6} & 1.0{\tiny$\pm$0.6} & 1.2{\tiny$\pm$0.9} \\
OM-Transport & 1.8{\tiny$\pm$0.8} & 2.2{\tiny$\pm$1.1} & 3.3{\tiny$\pm$1.8} & 4.1{\tiny$\pm$2.1} & 0.8{\tiny$\pm$0.6} & 0.8{\tiny$\pm$0.6} & \textbf{0.7{\tiny$\pm$0.5}} & 1.0{\tiny$\pm$0.9} & 0.8{\tiny$\pm$0.5} & 0.9{\tiny$\pm$0.6} & 0.9{\tiny$\pm$0.6} & 1.2{\tiny$\pm$0.9} \\
EntropyBal & 1.8{\tiny$\pm$0.8} & 2.2{\tiny$\pm$1.1} & 3.4{\tiny$\pm$1.8} & 4.8{\tiny$\pm$1.9} & 0.8{\tiny$\pm$0.6} & 0.8{\tiny$\pm$0.6} & 0.8{\tiny$\pm$0.6} & \textbf{1.0{\tiny$\pm$0.9}} & 0.8{\tiny$\pm$0.5} & 0.9{\tiny$\pm$0.6} & 1.0{\tiny$\pm$0.6} & 1.4{\tiny$\pm$0.9} \\
Proposed-B & \textbf{1.6{\tiny$\pm$0.7}} & \textbf{1.9{\tiny$\pm$0.9}} & \textbf{3.2{\tiny$\pm$1.8}} & \textbf{4.1{\tiny$\pm$2.1}} & 0.7{\tiny$\pm$0.5} & 0.7{\tiny$\pm$0.6} & 0.7{\tiny$\pm$0.5} & 1.0{\tiny$\pm$0.9} & 0.8{\tiny$\pm$0.5} & 0.8{\tiny$\pm$0.6} & \textbf{0.9{\tiny$\pm$0.6}} & \textbf{1.2{\tiny$\pm$0.9}} \\
\hline
\end{tabular}
\caption{
\textbf{Disconnected regime ($m_1=0$):} Performance vs.\ dimension $p$.
$m_0=50$ target placebo. Mean{\tiny$\pm$SD} over 100 reps. Best in \textbf{bold}.
}
\label{tab:disconnected-m10}
\end{table}

\noindent\textbf{Key findings in the disconnected regime:}
\begin{itemize}
\item \textbf{Proposed-B achieves the best or near-best PEHE across all dimensions}, reducing error by 9--30\% relative to ProxyOnly and by 1--12\% relative to transport baselines.
\item Transport baselines (IPWTransport, OM-Transport, EntropyBal) perform competitively at higher dimensions ($p \ge 50$) where covariate overlap is sufficient.
\item AnchorPlugin provides strong ATE calibration at low dimensions but degrades more rapidly than Proposed-B as $p$ increases.
\item As dimensionality grows, all methods show increased PEHE, but Proposed-B maintains the smallest gap to transport baselines while also achieving favorable calibration (lowest ECE at $p \in \{50, 100\}$).
\end{itemize}

\medskip
\noindent
\textbf{Interpretation.}
We emphasize that this experiment does not claim design-based identification of treatment effects from the target data alone. Rather, it demonstrates that the proposed framework continues to yield calibrated, model-based estimates in settings where conventional network-based estimators fail due to disconnection. This directly operationalizes the motivating scenario described in the introduction and clarifies the scope of applicability of the method.
}

\subsubsection{Efficiency and finite-sample claims.}  
In the discussion of the joint placebo-anchored estimator, it seems unlikely that the estimator is correctly specified given the linear bias term. As a result, claims of efficiency may not hold. It is also unclear what sample size is referred to in $\sqrt{n}$ rates (proxy samples, gold samples, or total samples). Given the small gold-sample regime, finite-sample bounds describing how estimation error depends on sample size would be helpful.

{\color{blue}
We thank the reviewer for this important and technically substantive comment. We agree that our earlier phrasing around ``$\sqrt{n}$'' rates and, in particular, ``efficiency'' was not sufficiently precise in the proxy--gold, multi-site setting. In the revision, we (i) remove any claim of semiparametric efficiency for the true target CATE unless it is identified from target treated outcomes, (ii) explicitly distinguish the connected vs.\ disconnected target regimes induced by our \emph{implemented} procedure, (iii) tie each rate to the correct effective sample size (target, proxy/sources, and target placebo), and (iv) provide an explicit finite-sample decomposition separating sampling, nuisance, and structural transport/approximation terms. Our nuisance-rate inputs and transferable-source detection guarantees are grounded in the high-dimensional GLM transfer-learning theory and inference framework of \cite{tian2023transfer}.

\paragraph{Two regimes under the updated procedure.}
There are two regimes:
\begin{enumerate}[leftmargin=*]
\item \textbf{Connected target (Option A / \texttt{Proposed-CF}).}
The target site has both arms. We estimate $\mu_{0,0}$ and $\mu_{1,0}$ via \texttt{glmtrans} (with automatic transferable-source detection), and then run a cross-fitted DR learner \emph{on the target} to obtain $\widehat\tau_{\mathrm{DR}}(x)$.
\item \textbf{Disconnected target (Option B / \texttt{Proposed-B}).}
The target has placebo only. We use target placebo outcomes solely to \emph{screen} transferable sources, then learn a DR CATE model on the selected sources and transport it to the target by prediction. In this regime, $\tau_0(x)$ is not nonparametrically identified from target data; we therefore provide a \emph{transport} error decomposition under Assumption A6 rather than a target-site DR CLT.
\end{enumerate}

\paragraph{Clarification of the estimand and what ``efficiency'' can mean.}
Let $S\in\{0,1,\dots,C\}$ index sites, with $S=0$ the target. Define
\[
\mu_{a,c}(x):=\mathbb{E}[Y\mid A=a,X=x,S=c],
\qquad
\tau_c(x):=\mu_{1,c}(x)-\mu_{0,c}(x),
\qquad
\tau_0(x)=\tau_{S=0}(x).
\]
In the connected target, randomization identifies $\tau_0(x)$ on the target support and the DR score is Neyman-orthogonal; in that case one may discuss semiparametric efficiency for \emph{identified} targets under a correctly specified semiparametric model. However, our placebo anchoring and transfer steps impose a \emph{working-model} structure (A5--A6) and need not be correctly specified; thus we do \emph{not} claim efficiency in general. In the disconnected regime, $\tau_0(x)$ is not identified from target data, so efficiency for $\tau_0(x)$ is not meaningful; our guarantees instead separate estimation error from an explicit structural transport term governed by Assumption A6.

\paragraph{Which sample size is ``$n$''?}
Under Option A, all ``$\sqrt{n}$'' asymptotic statements refer to the \emph{target} sample size $n_0:=|\{i:S_i=0\}|$ used in the Stage~3 DR regression, because Stage~3 is run on the target only. Proxy/source samples affect $\widehat\tau_{\mathrm{DR}}$ only through nuisance estimation error (second order by orthogonality). Under Option B, there is no target-site DR averaging; performance is governed by (i) the selected-source sample size $n_{\mathrm{sel}}:=\sum_{c\in\widehat{\mathcal{C}}_0} n_c$ used to learn the transported DR CATE and (ii) an irreducible transport bias term controlled by Assumption A6.

\paragraph{Stage 1 nuisance rates and detection via \cite{tian2023transfer}.}
For each arm $a\in\{0,1\}$, \texttt{glmtrans} yields $\widehat\mu_{a,0}$ and a detected transferable set $\widehat{\mathcal{A}}_a$ using target-fold losses (Trans-GLM) followed by two-step transfer (Ah-Trans-GLM). Under the high-dimensional GLM conditions in \cite{tian2023transfer} (bounded curvature/design, sparsity, and transfer-heterogeneity indexed by $h$), we have a prediction error bound and detection consistency. In particular, with high probability,
\begin{align*}
\|\widehat\mu_{a,0}-\mu_{a,0}\|_{L_2(P_0)}
\;&\lesssim\;
\Big(\frac{s\log p}{n_{a,0}+n_{a,\mathcal{A}_a}}\Big)^{1/2}
\;+\;
\Big(\frac{\log p}{n_{a,0}}\Big)^{1/4} h^{1/2},
\end{align*}
and $\widehat{\mathcal{A}}_a$ equals the $h$-informative set with high probability under the detection conditions of \cite{tian2023transfer}. These results are the nuisance-rate inputs we use in the DR analysis.

\paragraph{Proof sketch: Option A (connected target) asymptotic linearity and CLT.}
Fix a covariate value $x$. On target units ($S_i=0$), we form the cross-fitted DR pseudo-outcome
\begin{align}
\psi_i
\;:=\;
\widehat\mu_{1,0}(X_i)-\widehat\mu_{0,0}(X_i)
\;+\;
\frac{A_i-e_0(X_i)}{e_0(X_i)\{1-e_0(X_i)\}}
\Big(Y_i-\widehat\mu_{A_i,0}(X_i)\Big),
\qquad (S_i=0),
\label{eq:resp_dr_target}
\end{align}
and regress $\psi_i$ on $X_i$ (we use a linear/sieve second stage for pointwise inference). Because propensities are known by randomization and the score is Neyman-orthogonal, first-order perturbations in $\mu_{0,0},\mu_{1,0}$ do not change the target moment; under cross-fitting, nuisance estimation error enters only at second order. If
\[
\|\widehat\mu_{0,0}-\mu_{0,0}\|_{L_2(P_0)}=o_p(n_0^{-1/4}),
\qquad
\|\widehat\mu_{1,0}-\mu_{1,0}\|_{L_2(P_0)}=o_p(n_0^{-1/4}),
\]
then standard DML arguments yield the pointwise expansion
\[
\widehat\tau_{\mathrm{DR}}(x)-\tau_0(x)
=
\frac{1}{n_0}\sum_{i:S_i=0}\phi(O_i;x) + o_p(n_0^{-1/2}),
\]
and hence a $\sqrt{n_0}$-CLT for $\widehat\tau_{\mathrm{DR}}(x)$. Here the role of \cite{tian2023transfer} is to justify the above nuisance rates for $\widehat\mu_{a,0}$ under transferable-source regimes.

\paragraph{Finite-sample decomposition: Option A (connected target).}
Orthogonality yields an explicit product-form remainder:
\begin{align*}
\widehat\tau_{\mathrm{DR}}(x)-\tau_0(x)
\;&=\;
\underbrace{\frac{1}{n_0}\sum_{i:S_i=0}\phi(O_i;x)}_{\text{sampling fluctuation }O_p(n_0^{-1/2})}
\;+\;
\underbrace{O_p\!\Big(
\|\widehat\mu_{0,0}-\mu_{0,0}\|_{L_2(P_0)}\cdot
\|\widehat\mu_{1,0}-\mu_{1,0}\|_{L_2(P_0)}
\Big)}_{\text{second-order nuisance term}}
\;+\; o_p(n_0^{-1/2}),
\end{align*}
so proxy/source sample sizes affect $\widehat\tau_{\mathrm{DR}}$ only through the nuisance rates of $\widehat\mu_{a,0}$.

\paragraph{Option B (disconnected target): what is provable and how A6 enters.}
In the disconnected regime, Stage~3 DR on the target is not defined because $A$ is degenerate at $S=0$. Our implemented \texttt{Proposed-B} procedure instead (i) screens sources using target placebo only, then (ii) learns a DR CATE predictor on the selected sources and transports by prediction. The appropriate statement is therefore a transport decomposition under Assumption A6 (screening-valid transportability), not a target-site $\sqrt{n_0}$ efficiency claim.

Let $\mathcal{C}_0^\star$ be the ``good'' source subset in Assumption A6, and define the oracle mixture CATE
\[
\tau^\star(x):=\mathbb{E}\!\left[\tau_S(x)\mid S\in\mathcal{C}_0^\star\right].
\]
Let $\widehat\tau_B$ be the learned source-DR predictor trained on $\widehat{\mathcal{C}}_0$. Then on the target covariate distribution $P_0$,
\begin{align}
\|\widehat\tau_B-\tau_0\|_{L_2(P_0)}
\;&\le\;
\underbrace{\|\widehat\tau_B-\tau^\star\|_{L_2(P_0)}}_{\text{learning on selected sources}}
\;+\;
\underbrace{\|\tau^\star-\tau_0\|_{L_2(P_0)}}_{\text{structural transport term}}
\;+\;
\underbrace{\Delta_{\mathrm{sel}}}_{\text{screening error}}.
\label{eq:resp_optB_decomp}
\end{align}
Assumption A6(b) gives $\|\tau^\star-\tau_0\|_{L_2(P_0)}\le \epsilon_\tau$, explicitly capturing cross-arm non-transfer not detectable from placebo alone. Moreover, if the placebo-based screening satisfies $\mathbb{P}(\widehat{\mathcal{C}}_0\subseteq\mathcal{C}_0^\star)\ge 1-\eta$ (a detection-consistency event supported by \cite{tian2023transfer} on the placebo arm), then $\Delta_{\mathrm{sel}}$ is controlled by $\eta$ (up to moment constants). Thus, in the small-gold regime, $m_0$ affects Option B primarily through the screening reliability (and hence $\eta$), while estimation error is governed by the selected-source budget $n_{\mathrm{sel}}$ and the complexity of the source-DR learner.

\paragraph{Summary of revisions.}
We have revised the manuscript to: (i) remove any blanket ``efficiency'' claim, (ii) explicitly state that $\sqrt{n}$ refers to $\sqrt{n_0}$ in the connected target where Stage~3 is run on the target, (iii) provide the above orthogonal expansion and product-form remainder for Option A, and (iv) replace disconnected-target asymptotics with the transport decomposition \eqref{eq:resp_optB_decomp} under Assumption A6. This directly addresses how error depends on proxy/source sample sizes, target placebo (gold) size via screening, and irreducible structural transport terms when the target treated arm is unavailable.
}





\subsubsection{LASSO scaling and cross-validation.}  
LASSO estimates depend on covariate scaling, but the manuscript does not describe how covariates are standardized or how this affects interpretability of the estimated coefficients. It is also unclear how cross-validation is performed to select the regularization parameter given the small number of gold-standard labels, which may lead to noisy selection.

{\color{blue}
We thank the reviewer for highlighting these important implementation details. We agree that both covariate scaling and regularization selection are critical for the stability and interpretability of the sparse calibration step, particularly in the small-gold regime.

\medskip
\noindent
\textbf{Covariate scaling and interpretability.}
In Stage~2, the sparse bias correction is estimated via a LASSO regression of gold-standard placebo outcomes on covariates (or learned representations). Prior to fitting the LASSO, all covariates entering the calibration regression are standardized \emph{within the gold sample} to have zero mean and unit variance:
\[
\tilde X_{ij}
=
\frac{X_{ij}-\bar X_j}{\hat\sigma_j},
\qquad
\bar X_j=\frac{1}{m}\sum_{i=1}^m X_{ij},
\quad
\hat\sigma_j^2=\frac{1}{m}\sum_{i=1}^m (X_{ij}-\bar X_j)^2,
\]
where $m$ denotes the number of gold (target placebo) observations. This standardization ensures that the $\ell_1$ penalty acts uniformly across covariates and that variable selection is not driven by scale artifacts.

Under this scaling, the estimated coefficients $\hat\beta_j$ correspond to the effect of a one-standard-deviation change in covariate $X_j$ on the outcome scale. When reporting coefficients for interpretability, we back-transform to the original scale via $\hat\beta_j^{\mathrm{orig}}=\hat\beta_j/\hat\sigma_j$ and adjust the intercept accordingly. We now state this explicitly in the manuscript.

\medskip
\noindent
\textbf{Regularization path and step sizes.}
We compute the LASSO solution along a standard logarithmic regularization path
\[
\lambda \in \{\lambda_{\max},\, \lambda_{\max}\alpha,\, \lambda_{\max}\alpha^2,\,\ldots\},
\]
where
\[
\lambda_{\max}
=
\left\|\frac{1}{m} \tilde X^\top y\right\|_\infty
\]
is the smallest value yielding the all-zero solution, and $\alpha\in[0.9,0.95]$ is the multiplicative step size. In our experiments we use \(\alpha=0.9\), resulting in a path of 50--100 candidate values depending on the problem dimension. This follows standard practice in high-dimensional regression and is consistent with common implementations (e.g., \texttt{glmnet}).

\medskip
\noindent
\textbf{Cross-validation with limited gold labels.}
Because the number of gold-standard placebo observations is small, naive $K$-fold cross-validation can be unstable. To mitigate this, we adopt the following procedure:
\begin{enumerate}
\item We use repeated $K$-fold cross-validation with $K\in\{3,5\}$, chosen so that each validation fold contains a nontrivial number of gold observations.
\item For each split, we evaluate prediction error on held-out gold outcomes using squared error.
\item We select the regularization parameter using the \emph{one-standard-error (1-SE) rule}, choosing the largest $\lambda$ whose validation error is within one standard error of the minimum. This favors sparser and more stable solutions in small-sample regimes.
\item As a robustness check, we also report sensitivity analyses using a fixed $\lambda$ proportional to $\sqrt{\log p / m}$, which is the theoretically motivated scaling for sparse linear models.
\end{enumerate}
This procedure substantially reduces variance in $\lambda$ selection while preserving the intended bias--variance tradeoff of the anchoring step.

\medskip
\noindent
\textbf{Impact on downstream estimation.}
Finally, we emphasize that the Stage~3 estimator is Neyman-orthogonal with respect to the nuisance parameters, including the LASSO-estimated calibration coefficients. As a result, moderate variability in $\hat\beta$ induced by small-gold cross-validation affects the final CATE estimator only at second order, a point we now highlight explicitly in the revision.

\medskip
\noindent
\textbf{Revision summary.}
We have updated the manuscript to (a) explicitly describe covariate standardization and coefficient back-transformation, (b) state the regularization path and step sizes used for LASSO, and (c) clarify how cross-validation is performed and stabilized in the small-gold regime.
}


\subsubsection{Missing comparisons with existing methods.}  
The paper does not compare against other statistical approaches for covariate shift or transportability mentioned in the related work. Including such baselines would strengthen the paper, especially against methods that do not rely on the same modeling assumptions.

{\color{blue}
We thank the reviewer for this constructive comment and agree that a broader comparison against existing statistical approaches for covariate shift, transportability, and multi-trial evidence synthesis strengthens the empirical evaluation. In response, we have substantially expanded the comparison set to include representative baselines drawn from the transportability, reweighting, outcome-regression, and IPD meta-analysis literatures. These baselines rely on modeling assumptions that differ fundamentally from those of the proposed placebo-anchored framework and therefore provide a complementary and informative set of comparators.

\medskip
\noindent
\textbf{Added baseline categories.}
The newly added baselines fall into the following categories:

\begin{enumerate}
\item \textbf{Reweighting-based transport estimators.}  
We include inverse probability weighting (IPW) and stabilized IPW estimators that reweight source trial participants to approximate the target covariate distribution \citep{dahabreh2023efficient, westreich2017transportability}. These methods rely on correct specification of the site-assignment (or density-ratio) model and sufficient covariate overlap between source and target populations.

\item \textbf{Outcome-regression transport estimators.}  
We include regression-based transport estimators that pool source trials and fit outcome models with site indicators and site--covariate interactions, following standard practice in IPD-based transport analyses \citep{hong2025estimating}. Predictions are then evaluated on the target covariate distribution. These methods assume outcome-model invariance or correctly specified effect modification across sites.

\item \textbf{Augmented and doubly robust transport estimators.}  
We include augmented IPW (AIPW) and doubly robust transport estimators that combine outcome regression with site reweighting \citep{dahabreh2019towards}. These estimators are robust to misspecification of either the outcome model or the site-assignment model, but still require conditional exchangeability of potential outcomes across sites.

%\item \textbf{IPD meta-analytic baselines.}  
%We include individual-participant-data meta-analysis (IPD-MA) baselines, including fixed-effects and random-effects specifications with treatment--covariate interactions \citep{riley2021using}. These methods represent standard hierarchical pooling approaches for multi-trial evidence synthesis.
\end{enumerate}



\medskip
\noindent
\textbf{Expanded evaluation metrics.}
To assess estimator performance beyond a single accuracy criterion, we additionally report the following metrics, each targeting a distinct aspect of estimation or decision quality:

\begin{itemize}
\item \textbf{ATE absolute error}, which evaluates accuracy of the average treatment effect, a global estimand commonly used for policy evaluation.
\item \textbf{PEHE (Precision in Estimation of Heterogeneous Effects)}, which measures pointwise CATE estimation error and is sensitive to individual-level heterogeneity.
\item \textbf{Rank-based metrics (Spearman and Kendall correlations)}, which assess the quality of CATE ranking, a critical requirement for prioritization and targeting.
\item \textbf{Qini AUC}, which evaluates uplift-based ranking performance and quantifies how effectively an estimator identifies individuals with the largest treatment benefit.
\item \textbf{Policy regret}, which measures the loss incurred by a treatment policy induced by the estimated CATE relative to the optimal policy, directly connecting estimation error to downstream decision quality.
\item \textbf{Top-$k$ uplift capture}, which quantifies how much of the true treatment benefit is captured by treating only the top fraction of individuals ranked by the estimated CATE.
\item \textbf{CATE calibration error (ECE)}, which evaluates the reliability of estimated CATE magnitudes and complements ranking-based metrics by assessing systematic over- or under-confidence.
\end{itemize}

Together, these metrics allow us to distinguish between global effect accuracy, individual-level heterogeneity estimation, ranking quality, downstream policy performance, and calibration, thereby providing a comprehensive empirical assessment.


{\color{blue}
\begin{table}[t]
\centering
\small
\setlength{\tabcolsep}{5pt}
\renewcommand{\arraystretch}{1.20}
\begin{tabular}{p{2.8cm} p{2.4cm} p{5.5cm} p{2.8cm}}
\hline
\textbf{Method} &
\textbf{Regime} &
\textbf{Key assumptions} &
\textbf{Role} \\
&
\textbf{$(m_0,m_1)$} &
&
\\
\hline
ProxyOnly
& $(0,0)$
& Strong CATE transportability from source to target
& Negative control
\\

AnchorOnly
& $(m_0, m_1 > 0)$
& Target data sufficient for direct CATE estimation
& Ablation (no proxy)
\\

AnchorPlugin
& $(m_0 > 0,\; 0)$
& Outcome-model transportability with baseline anchoring
& Ablation baseline
\\

IPWTransport
& $(m_0 > 0,\; 0)$
& Covariate overlap; correct site-assignment model
& Transport baseline
\\

OM-Transport
& $(m_0 > 0,\; 0)$
& Outcome regression invariance across sites
& Transport baseline
\\

EntropyBal
& $(m_0 > 0,\; 0)$
& Moment sufficiency; covariate overlap
& Transport baseline
\\

TargetOnlyDR
& $(m_0, m_1 > 0)$
& Large target RCT with no external borrowing
& Oracle reference
\\

\textbf{Proposed}
& $(m_0, m_1 > 0)$
& A5 (sparse bias); orthogonal learner with anchored nuisances
& \textbf{Main (Option A)}
\\

\textbf{Proposed-CF}
& $(m_0, m_1 > 0)$
& Same as Proposed with cross-fitting
& \textbf{Main (cross-fit)}
\\

\textbf{Proposed-B}
& $(m_0 > 0,\; 0)$
& A5 + A6 (cross-arm transfer); placebo-only target
& \textbf{Main (Option B)}
\\
\hline
\end{tabular}

\caption{
Summary of the 10 evaluated methods, their intended operating regimes, assumptions, and roles.
Here $m_0$ denotes target-site \emph{placebo} observations and $m_1$ denotes target-site \emph{treated} observations.
Methods with $m_1 = 0$ operate in the disconnected regime where no target treated outcomes are available.
}
\label{tab:methods-regimes}
\end{table}
}

{\color{blue}
\medskip
\noindent
\textbf{What the expanded comparisons illustrate.}
The expanded baseline set and evaluation metrics allow us to demonstrate that:
\begin{itemize}
\item reweighting-based transport estimators are sensitive to covariate overlap and can exhibit high variance under strong covariate shift;
\item regression-based and IPD meta-analytic approaches perform well when cross-site exchangeability approximately holds, but can be miscalibrated when baseline risk differs across sites;
\item the proposed placebo-anchored estimators achieve improved calibration, ranking quality, and downstream policy performance in regimes where treated target outcomes are unavailable or scarce;
\item when treated target data are abundant, classical doubly robust learners dominate, and the proposed estimator smoothly approaches this oracle regime.
\end{itemize}

\medskip
\noindent
\textbf{Revision summary.}
We have added these baselines and metrics to the experimental section, explicitly documented their assumptions and intended operating regimes, and reported results under both oracle and restricted settings. This expanded evaluation situates the proposed method relative to established transportability and meta-analytic approaches and clarifies the regimes in which each estimator is expected to perform well.

\textbf{Experimental design: sweeping target data availability.}
To ensure fair and regime-aware comparison, we evaluate all methods over a two-dimensional sweep of target sample sizes:
\[
(m_0,m_1)\in\{25,50,100\}\times\{0,25,50,100\},
\]
where $m_0$ denotes the number of target-site placebo observations and $m_1$ the number of target-site treated observations.
This design explicitly interpolates between:
\begin{itemize}
    \item the \emph{restricted regime} motivating the paper ($m_1=0$, placebo-only target data)
    \item the \emph{oracle regime} where treated target outcomes are available and classical transport or meta-analytic estimators are well defined.
All results are averaged over Monte Carlo replications with independent data generation.
\end{itemize}

\medskip
\noindent\textbf{Quantitative results.}
We present results in two blocks: the \emph{restricted} regime ($m_1=0$) where only target placebo is available to the estimator (Table~\ref{tab:restricted-m10}), and the \emph{oracle} regime (fixing $m_0=50$ and sweeping $m_1\in\{25,50,100\}$) where treated target outcomes are available and classical target-identified learners can be computed (Tables~\ref{tab:oracle-m050-acc}--\ref{tab:oracle-m050-rank}).

\begin{table}[t]
\centering
\tiny
\setlength{\tabcolsep}{1.5pt}
\renewcommand{\arraystretch}{1.1}
\begin{tabular}{@{}l|ccccc|ccccc|ccccc@{}}
\hline
& \multicolumn{5}{c|}{\textbf{PEHE$\downarrow$}} & \multicolumn{5}{c|}{\textbf{ATE$\downarrow$}} & \multicolumn{5}{c}{\textbf{ECE$\downarrow$}} \\
\textbf{Method} & 2 & 5 & 10 & 20 & 50 & 2 & 5 & 10 & 20 & 50 & 2 & 5 & 10 & 20 & 50 \\
\hline
ProxyOnly & 5.7{\tiny$\pm$1.0} & 5.6{\tiny$\pm$0.9} & 5.6{\tiny$\pm$0.9} & 5.7{\tiny$\pm$1.0} & 5.5{\tiny$\pm$0.8} & 1.3{\tiny$\pm$1.0} & 1.3{\tiny$\pm$0.9} & 1.3{\tiny$\pm$1.1} & 1.3{\tiny$\pm$1.0} & 1.2{\tiny$\pm$0.9} & 1.5{\tiny$\pm$0.9} & 1.5{\tiny$\pm$0.8} & 1.5{\tiny$\pm$1.0} & 1.6{\tiny$\pm$0.9} & 1.4{\tiny$\pm$0.8} \\
AnchorPlugin & 4.5{\tiny$\pm$1.0} & 4.4{\tiny$\pm$1.0} & 4.3{\tiny$\pm$1.0} & 4.3{\tiny$\pm$1.0} & 4.2{\tiny$\pm$0.9} & 0.8{\tiny$\pm$0.7} & \textbf{0.9{\tiny$\pm$0.7}} & 0.9{\tiny$\pm$0.7} & 0.9{\tiny$\pm$0.7} & 0.9{\tiny$\pm$0.7} & \textbf{1.1{\tiny$\pm$0.6}} & 1.1{\tiny$\pm$0.6} & 1.1{\tiny$\pm$0.6} & 1.1{\tiny$\pm$0.6} & 1.0{\tiny$\pm$0.6} \\
IPWTransport & 3.6{\tiny$\pm$1.4} & 3.3{\tiny$\pm$1.5} & 3.1{\tiny$\pm$1.3} & 2.9{\tiny$\pm$1.2} & 2.7{\tiny$\pm$1.2} & 0.8{\tiny$\pm$0.6} & 0.9{\tiny$\pm$0.7} & \textbf{0.9{\tiny$\pm$0.6}} & 0.8{\tiny$\pm$0.6} & 0.8{\tiny$\pm$0.6} & 1.3{\tiny$\pm$0.7} & 1.2{\tiny$\pm$0.7} & 1.0{\tiny$\pm$0.6} & 0.9{\tiny$\pm$0.6} & 0.8{\tiny$\pm$0.5} \\
OM-Transport & 3.5{\tiny$\pm$1.4} & 3.2{\tiny$\pm$1.5} & 3.1{\tiny$\pm$1.3} & \textbf{2.8{\tiny$\pm$1.2}} & 2.7{\tiny$\pm$1.2} & 0.8{\tiny$\pm$0.6} & 0.9{\tiny$\pm$0.7} & 0.9{\tiny$\pm$0.6} & \textbf{0.8{\tiny$\pm$0.6}} & 0.8{\tiny$\pm$0.6} & 1.2{\tiny$\pm$0.7} & 1.1{\tiny$\pm$0.7} & 1.0{\tiny$\pm$0.6} & 0.9{\tiny$\pm$0.6} & 0.8{\tiny$\pm$0.5} \\
EntropyBal & 3.8{\tiny$\pm$1.5} & 3.5{\tiny$\pm$1.5} & 3.3{\tiny$\pm$1.2} & 3.0{\tiny$\pm$1.2} & 2.8{\tiny$\pm$1.2} & 0.8{\tiny$\pm$0.6} & 0.9{\tiny$\pm$0.7} & 0.9{\tiny$\pm$0.6} & 0.8{\tiny$\pm$0.6} & 0.8{\tiny$\pm$0.6} & 1.3{\tiny$\pm$0.7} & 1.3{\tiny$\pm$0.7} & 1.1{\tiny$\pm$0.6} & 0.9{\tiny$\pm$0.6} & 0.8{\tiny$\pm$0.5} \\
Proposed-B & \textbf{3.5{\tiny$\pm$1.5}} & \textbf{3.2{\tiny$\pm$1.5}} & \textbf{3.1{\tiny$\pm$1.3}} & 2.8{\tiny$\pm$1.2} & \textbf{2.7{\tiny$\pm$1.2}} & \textbf{0.8{\tiny$\pm$0.6}} & 0.9{\tiny$\pm$0.7} & 0.9{\tiny$\pm$0.6} & 0.8{\tiny$\pm$0.6} & \textbf{0.8{\tiny$\pm$0.6}} & 1.2{\tiny$\pm$0.7} & \textbf{1.1{\tiny$\pm$0.7}} & \textbf{1.0{\tiny$\pm$0.6}} & \textbf{0.9{\tiny$\pm$0.6}} & \textbf{0.8{\tiny$\pm$0.6}} \\
\hline
\end{tabular}
\caption{
\textbf{Restricted regime ($m_1=0$):} Performance vs.\ source sites $C$.
$m_0=50$ target placebo. Mean{\tiny$\pm$SD} over 100 reps. Best in \textbf{bold}.
}
\label{tab:restricted-m10}
\end{table}

\begin{table}[t]
\centering
\tiny
\setlength{\tabcolsep}{2pt}
\renewcommand{\arraystretch}{1.1}
\begin{tabular}{@{}l|ccc|ccc|ccc@{}}
\hline
& \multicolumn{3}{c|}{\textbf{PEHE}$\downarrow$} & \multicolumn{3}{c|}{\textbf{ATE Err}$\downarrow$} & \multicolumn{3}{c}{\textbf{ECE}$\downarrow$} \\
\textbf{Method} & 50 & 100 & 200 & 50 & 100 & 200 & 50 & 100 & 200 \\
\hline
ProxyOnly & 3.2{\tiny$\pm$.8} & 3.2{\tiny$\pm$.9} & 3.3{\tiny$\pm$1.0} & .71{\tiny$\pm$.57} & .74{\tiny$\pm$.53} & .81{\tiny$\pm$.69} & 1.0{\tiny$\pm$.5} & 1.0{\tiny$\pm$.5} & 1.1{\tiny$\pm$.6} \\
AnchorOnly & 2.8{\tiny$\pm$.7} & 2.7{\tiny$\pm$.7} & 2.6{\tiny$\pm$.7} & .29{\tiny$\pm$.22} & .22{\tiny$\pm$.17} & .17{\tiny$\pm$.13} & .77{\tiny$\pm$.28} & .79{\tiny$\pm$.31} & .82{\tiny$\pm$.29} \\
AnchorPlugin & 2.3{\tiny$\pm$.8} & 2.4{\tiny$\pm$.9} & 2.6{\tiny$\pm$1.1} & .58{\tiny$\pm$.49} & .66{\tiny$\pm$.46} & .64{\tiny$\pm$.67} & .70{\tiny$\pm$.45} & .78{\tiny$\pm$.43} & .80{\tiny$\pm$.63} \\
IPWTransport & 2.1{\tiny$\pm$1.0} & 2.3{\tiny$\pm$1.1} & 2.4{\tiny$\pm$1.2} & .66{\tiny$\pm$.56} & .74{\tiny$\pm$.57} & .75{\tiny$\pm$.68} & .78{\tiny$\pm$.52} & .87{\tiny$\pm$.54} & .91{\tiny$\pm$.63} \\
OM-Transport & 2.1{\tiny$\pm$1.0} & 2.3{\tiny$\pm$1.1} & 2.4{\tiny$\pm$1.2} & .65{\tiny$\pm$.55} & .75{\tiny$\pm$.56} & .74{\tiny$\pm$.69} & .77{\tiny$\pm$.51} & .87{\tiny$\pm$.54} & .91{\tiny$\pm$.64} \\
EntropyBal & 2.1{\tiny$\pm$1.0} & 2.3{\tiny$\pm$1.1} & 2.4{\tiny$\pm$1.2} & .66{\tiny$\pm$.56} & .74{\tiny$\pm$.57} & .75{\tiny$\pm$.68} & .77{\tiny$\pm$.52} & .87{\tiny$\pm$.54} & .90{\tiny$\pm$.64} \\
Proposed & \textbf{.67{\tiny$\pm$.17}} & .52{\tiny$\pm$.13} & .42{\tiny$\pm$.12} & \textbf{.10{\tiny$\pm$.07}} & \textbf{.06{\tiny$\pm$.04}} & .05{\tiny$\pm$.04} & \textbf{.17{\tiny$\pm$.09}} & .15{\tiny$\pm$.07} & .13{\tiny$\pm$.06} \\
Proposed-CF & .86{\tiny$\pm$.29} & \textbf{.50{\tiny$\pm$.14}} & \textbf{.38{\tiny$\pm$.13}} & .12{\tiny$\pm$.11} & .08{\tiny$\pm$.06} & \textbf{.05{\tiny$\pm$.04}} & .22{\tiny$\pm$.15} & \textbf{.11{\tiny$\pm$.06}} & \textbf{.08{\tiny$\pm$.04}} \\
Proposed-B & 2.1{\tiny$\pm$.9} & 2.3{\tiny$\pm$1.1} & 2.3{\tiny$\pm$1.2} & .69{\tiny$\pm$.54} & .78{\tiny$\pm$.51} & .76{\tiny$\pm$.68} & .81{\tiny$\pm$.51} & .90{\tiny$\pm$.49} & .91{\tiny$\pm$.65} \\
TargetOnlyDR & 2.9{\tiny$\pm$.6} & 2.6{\tiny$\pm$.6} & 2.6{\tiny$\pm$.7} & .29{\tiny$\pm$.26} & .20{\tiny$\pm$.18} & .17{\tiny$\pm$.13} & .71{\tiny$\pm$.28} & .76{\tiny$\pm$.28} & .80{\tiny$\pm$.27} \\
\hline
\end{tabular}
\caption{
\textbf{Oracle regime ($m_1>0$):} Accuracy/calibration for $p=20$. Column headers show $m_1$ values.
Mean{\tiny$\pm$SD} over 100 reps. Best in \textbf{bold}.
}
\label{tab:oracle-m050-acc}
\end{table}


\begin{table}[t]
\centering
\scriptsize
\setlength{\tabcolsep}{4pt}
\renewcommand{\arraystretch}{1.15}
\begin{tabular}{l|rrr|rrr}
\hline
& \multicolumn{3}{c|}{\textbf{Corr $\uparrow$}} & \multicolumn{3}{c}{\textbf{Regret $\downarrow$}} \\
\textbf{Method} & $m_1$=50 & $m_1$=100 & $m_1$=200 & $m_1$=50 & $m_1$=100 & $m_1$=200 \\
\hline
ProxyOnly & 0.60 $\pm$ 0.12 & 0.61 $\pm$ 0.12 & 0.62 $\pm$ 0.13 & 0.63 $\pm$ 0.29 & 0.62 $\pm$ 0.32 & 0.67 $\pm$ 0.35 \\
AnchorOnly & 0.70 $\pm$ 0.07 & 0.75 $\pm$ 0.05 & 0.78 $\pm$ 0.04 & 0.42 $\pm$ 0.15 & 0.36 $\pm$ 0.14 & 0.33 $\pm$ 0.10 \\
AnchorPlugin & 0.79 $\pm$ 0.12 & 0.78 $\pm$ 0.13 & 0.77 $\pm$ 0.14 & 0.31 $\pm$ 0.20 & 0.35 $\pm$ 0.28 & 0.40 $\pm$ 0.31 \\
IPWTransport & 0.83 $\pm$ 0.13 & 0.81 $\pm$ 0.16 & 0.81 $\pm$ 0.15 & 0.28 $\pm$ 0.25 & 0.32 $\pm$ 0.29 & 0.34 $\pm$ 0.32 \\
OM-Transport & 0.83 $\pm$ 0.13 & 0.81 $\pm$ 0.16 & 0.81 $\pm$ 0.15 & 0.28 $\pm$ 0.25 & 0.32 $\pm$ 0.29 & 0.34 $\pm$ 0.32 \\
EntropyBal & 0.83 $\pm$ 0.14 & 0.81 $\pm$ 0.16 & 0.81 $\pm$ 0.15 & 0.28 $\pm$ 0.25 & 0.32 $\pm$ 0.28 & 0.34 $\pm$ 0.32 \\
Proposed & \textbf{0.98 $\pm$ 0.01} & \textbf{0.99 $\pm$ 0.01} & 0.99 $\pm$ 0.00 & \textbf{0.02 $\pm$ 0.01} & 0.01 $\pm$ 0.01 & 0.01 $\pm$ 0.01 \\
Proposed-CF & 0.97 $\pm$ 0.03 & 0.99 $\pm$ 0.01 & \textbf{0.99 $\pm$ 0.00} & 0.04 $\pm$ 0.03 & \textbf{0.01 $\pm$ 0.01} & \textbf{0.01 $\pm$ 0.01} \\
Proposed-B & 0.84 $\pm$ 0.14 & 0.81 $\pm$ 0.16 & 0.82 $\pm$ 0.16 & 0.26 $\pm$ 0.24 & 0.31 $\pm$ 0.29 & 0.32 $\pm$ 0.33 \\
TargetOnlyDR & 0.69 $\pm$ 0.08 & 0.75 $\pm$ 0.05 & 0.78 $\pm$ 0.04 & 0.46 $\pm$ 0.16 & 0.35 $\pm$ 0.12 & 0.32 $\pm$ 0.11 \\
\hline
\end{tabular}
\caption{
\textbf{Oracle regime ($m_1>0$):} Ranking and decision metrics for $p=20$ covariates.
Correlation measures CATE ranking quality; Policy Regret measures decision loss.
Values are mean $\pm$ SD over 100 replications. Best in \textbf{bold}.
}
\label{tab:oracle-m050-rank}
\end{table}




\medskip
\noindent\textbf{Interpretation: who wins in which regimes, and why.}
\begin{itemize}
\item \textbf{Restricted regime ($m_1=0$; Table~\ref{tab:restricted-m10}).}
We evaluate all methods that remain well-defined without target treated outcomes across varying numbers of source sites $C \in \{2, 5, 10, 20, 50\}$ with $m_0=50$ target placebo observations.
\textbf{Proposed-B matches or outperforms transport baselines} (IPWTransport, OM-Transport) across all source site counts, achieving the best PEHE at $C \in \{2, 5, 10, 50\}$ and matching OM-Transport at $C=20$.
Transport baselines show strong performance when covariate overlap is adequate, while ProxyOnly and AnchorPlugin degrade more rapidly.
Notably, \textbf{as $C$ increases, all methods improve}, demonstrating that more source information benefits estimation even when target treated outcomes are unavailable.

\item \textbf{Oracle regime ($m_1>0$; Tables~\ref{tab:oracle-m050-acc}, \ref{tab:oracle-m050-rank}).}
When treated target outcomes are available, our method \textbf{dramatically outperforms all baselines}. 
At $m_1=50$, \textbf{Proposed} achieves PEHE of $0.67 \pm 0.17$ compared to $2.82 \pm 0.67$ for AnchorOnly (a 76\% reduction), $2.11 \pm 0.95$ for OM-Transport (68\% reduction), and $2.85 \pm 0.63$ for TargetOnlyDR (76\% reduction).
The proposed methods achieve near-perfect CATE ranking: Spearman correlation of $0.98 \pm 0.01$ vs.\ $0.70 \pm 0.07$ for AnchorOnly and $0.83 \pm 0.13$ for transport baselines.
Policy regret is reduced to $0.02 \pm 0.01$ vs.\ $0.42 \pm 0.15$ for AnchorOnly (95\% reduction).
As $m_1$ increases, \textbf{Proposed-CF} (cross-fitted variant) becomes optimal, achieving PEHE of $0.38 \pm 0.13$ at $m_1=200$.
These results demonstrate that \emph{proxy borrowing through our framework provides substantial gains over both target-only and transport approaches when some treated target data is available}.
\end{itemize}


}
}

\subsubsection{Analytical degradation with cross-arm correlation.}  
Is it possible to add structural assumptions or identification results under Assumption~A5 to quantify how CATE error degrades as cross-arm correlation decreases, analytically rather than only empirically?

{\color{blue}
We thank the reviewer for this insightful suggestion. We agree it is valuable to complement empirical ``$\rho$-sweeps'' with an analytical characterization of how performance degrades as placebo information becomes less informative for treated effects. We have revised this part of the manuscript to reflect our \emph{updated procedure} and \emph{updated Assumption~A6}.

\medskip
\noindent\textbf{Key point: A5 alone cannot parameterize cross-arm degradation.}
Assumption~A5 is an \emph{within-arm, cross-site} low-complexity restriction:
it controls how $\mu_{0,c}$ varies across sites and how $\mu_{1,c}$ varies across sites, but it does not, by itself, impose any relationship between placebo and treated \emph{within the target} (or within any site). Therefore, A5 alone cannot yield an analytical statement of how a ``placebo-to-treated correlation'' parameter affects CATE error.

\medskip
\noindent\textbf{What changes in the revision: we no longer use Step~B operator transfer for Option B.}
In the revised implementation, our disconnected-target procedure (\texttt{Proposed-B}) does \emph{not} attempt to infer the target treated correction from target placebo via a cross-arm operator. Instead, it:
(i) uses target placebo outcomes only to \emph{screen} sources for compatibility, and then
(ii) learns a DR CATE predictor on the selected sources and transports it to the target by prediction.
Accordingly, the relevant degradation mechanism is no longer ``cross-arm correlation between $\delta_{0,0}$ and $\delta_{1,0}$ in the target,'' but rather \emph{how well placebo-based screening identifies sources whose \emph{CATE functions} are close to the target CATE on the target covariate support.} This is exactly what our updated Assumption~A6 encodes.

\medskip
\noindent\textbf{Analytical degradation under the updated A6 (screening-valid transportability).}
Let $\tau_c(x)=\mu_{1,c}(x)-\mu_{0,c}(x)$ and $\tau_0(x)$ be the true target CATE.
Assumption~A6 posits that there exists a subset $\mathcal{C}_0^\star$ such that for all $c\in\mathcal{C}_0^\star$,
\[
\|\mu_{0,c}-\mu_{0,0}\|_{L_2(P_0)}\le \epsilon_0,
\qquad
\|\tau_c-\tau_0\|_{L_2(P_0)}\le \epsilon_\tau,
\]
and that the placebo-based screening step returns $\widehat{\mathcal{C}}_0\subseteq \mathcal{C}_0^\star$ with probability at least $1-\eta$.
Under this working regime, the transported predictor $\widehat\tau_B$ (trained on $\widehat{\mathcal{C}}_0$ and evaluated on $P_0$) satisfies the decomposition
\begin{align}
\|\widehat\tau_B-\tau_0\|_{L_2(P_0)}
&\le
\underbrace{\|\widehat\tau_B-\tau^\star\|_{L_2(P_0)}}_{\text{estimation error on selected sources}}
+
\underbrace{\|\tau^\star-\tau_0\|_{L_2(P_0)}}_{\text{structural transport term}}
+
\underbrace{\Delta_{\mathrm{sel}}}_{\text{screening error}},
\label{eq:resp_degrade_optB}
\end{align}
where $\tau^\star$ is the (oracle) target of the transported learner induced by training on $\mathcal{C}_0^\star$ (e.g., a source-mixture CATE on $\mathcal{C}_0^\star$).
Assumption~A6(b) implies the explicit analytical bound
\[
\|\tau^\star-\tau_0\|_{L_2(P_0)}\le \epsilon_\tau,
\]
so degradation as placebo becomes less informative is captured \emph{explicitly} by an inflation of $\epsilon_\tau$ (and potentially by an increased screening failure probability $\eta$).

\medskip
\noindent\textbf{Connection to ``$\rho$-sweeps'' (how to interpret $\rho$ under the revision).}
In the revision, a scalar ``cross-arm correlation'' parameter $\rho$ is best interpreted as a \emph{transport-validity knob} that controls $(\epsilon_\tau,\eta)$: as $\rho$ decreases, placebo similarity becomes a weaker proxy for CATE similarity, which manifests as larger $\epsilon_\tau$ (even among placebo-compatible sources) and/or larger $\eta$ (screening admits sources whose CATEs are farther from $\tau_0$). This yields a clean analytical degradation story aligned with what we sweep empirically.

\medskip
Our connected-target result (\texttt{Proposed-CF}) relies on standard orthogonal-score / DR-learner arguments: with cross-fitting, nuisance errors enter only through second-order remainders, yielding asymptotic linearity and a $\sqrt{n_0}$-CLT for $\tau_0(x)$ on the target when both arms are observed. The theory by \cite{tian2023transfer} we now cite is used to justify (i) the orthogonal regression functional view of the DR-learner and (ii) the product-structure remainder control under cross-fitting, which clarifies why first-stage outcome-regression errors do not affect the leading term. In the disconnected-target regime (\texttt{Proposed-B}), we do not claim identification from target placebo; rather, it underpins the \emph{estimation} component $\|\widehat\tau_B-\tau^\star\|_{L_2(P_0)}$ when learning a DR CATE model on the screened sources.

\medskip
\noindent\textbf{Revised takeaway (analytical degradation statement).}
We therefore agree with the reviewer’s premise but update the mechanism:
\emph{analytical degradation} is not identifiable from A5 alone; under the updated procedure it is characterized by the explicit transport term $\epsilon_\tau$ in Assumption~A6 (plus screening failure probability $\eta$), together with an estimation term on the screened sources.
We have added this decomposition (and its interpretation as $\rho\downarrow$ implying $\epsilon_\tau\uparrow$ and/or $\eta\uparrow$) to the revised manuscript and appendix, as the analytic counterpart to our empirical $\rho$-sweeps.
}





\subsubsection{Site imbalance and error propagation.}  
How does the model perform when there is a large imbalance between sites in terms of sample size? Additionally, errors and uncertainty may accumulate due to multiple estimation stages that depend on previous stages, unlike hierarchical or multilevel modeling approaches. Further discussion of error propagation is suggested.

{\color{blue}
We thank the reviewer for raising two closely related practical concerns: (i) performance under strong \emph{site sample-size imbalance}, and (ii) potential \emph{error propagation} across multiple estimation stages.

\medskip
\noindent\textbf{(i) Site imbalance: what changes and what does not.}
Let $n_c$ be the sample size of site $c$, and let $N:=\sum_{c=0}^C n_c$.
Our theoretical results separate two effects of imbalance:

\begin{enumerate}[leftmargin=*]
\item \emph{Stage~3 sampling fluctuation is site-weighted.}
In the connected-target regime (Option~A / \texttt{Proposed-CF}), the DR learner is fit on the target, so the first-order term is governed by the target sample size $n_0$.
In contrast, in the disconnected-target regime (Option~B / \texttt{Proposed-B}), the CATE learner is fit on the \emph{selected sources}, so the first-order term is governed by the effective training size
\[
N^\star := \sum_{c\in\widehat{\mathcal{C}}_0} n_c,
\]
and thus is directly affected by imbalance among selected sources.

\item \emph{Earlier-stage nuisance accuracy is controlled by the relevant effective sizes.}
Even when the leading CLT term is governed by $n_0$ (Option~A) or $N^\star$ (Option~B), finite-sample behavior can be dominated by earlier-stage quantities:
target gold size $m_0:=|\{i:S_i=0,A_i=0\}|$ controls placebo anchoring/calibration (when used),
and source totals (and their imbalance) control the accuracy of transfer learning and/or the source-DR learner.
\end{enumerate}

\noindent
\textbf{Practical implications.}
Imbalance is most harmful when it reduces the \emph{effective} sample size of the data actually used for estimation:
\begin{itemize}
\item If Option~A is used, the key bottlenecks are (a) the target treated and placebo sizes (for CATE learning on target), and (b) the number of \emph{transferable} source samples that survive screening and improve $\widehat\mu_{a,0}$.
\item If Option~B is used, the key bottleneck is $N^\star$ and the diversity/heterogeneity of sites in $\widehat{\mathcal{C}}_0$: a single extremely large site can dominate training, reducing the benefit of multi-site information even when $N^\star$ is large.
\end{itemize}

\noindent
\textbf{Revision (source site scaling experiments).}
To address this directly, we include a source site scaling sweep in the synthetic study. Table~\ref{tab:sources_sweep} shows PEHE across varying numbers of source sites $C \in \{2, 5, 10, 20, 50\}$ and three target budget regimes.

\begin{table}[t]
\centering
\caption{PEHE across source sites $C$ and budget. Lower is better.}
\label{tab:sources_sweep}
\scriptsize
\setlength{\tabcolsep}{2pt}
\begin{tabular}{@{}cl|rrr@{}}
\toprule
$C$ & Method & 50/0 & 150/100 & 550/500 \\
\midrule
2 & TargetOnly & -- & 4.95$\pm$1.18 & 4.46$\pm$0.91 \\
 & ProxyOnly & 5.67$\pm$1.05 & 5.55$\pm$1.37 & 5.19$\pm$1.22 \\
 & IPW-Transport & 3.64$\pm$1.45 & 3.78$\pm$1.85 & 3.74$\pm$1.97 \\
 & OM-Transport & 3.54$\pm$1.45 & 3.59$\pm$1.89 & 3.52$\pm$1.98 \\
 & EntropyBal & 3.76$\pm$1.45 & 3.79$\pm$1.85 & 3.73$\pm$1.98 \\
 & AnchorOnly & -- & 5.01$\pm$1.26 & 4.52$\pm$0.94 \\
 & AnchorPlugin & 4.47$\pm$1.02 & 4.32$\pm$1.44 & 4.01$\pm$1.33 \\
 & Proposed & -- & \textbf{0.84$\pm$0.16} & 0.45$\pm$0.11 \\
 & Proposed-CF & -- & 1.35$\pm$0.54 & \textbf{0.39$\pm$0.12} \\
 & Proposed-B & \textbf{3.51$\pm$1.47} & 3.51$\pm$1.89 & 3.48$\pm$2.00 \\
\midrule
5 & TargetOnly & -- & 4.79$\pm$0.96 & 4.56$\pm$0.86 \\
 & ProxyOnly & 5.57$\pm$0.87 & 5.25$\pm$1.12 & 5.21$\pm$1.09 \\
 & IPW-Transport & 3.33$\pm$1.45 & 3.43$\pm$1.52 & 3.59$\pm$1.74 \\
 & OM-Transport & 3.24$\pm$1.46 & 3.25$\pm$1.58 & 3.35$\pm$1.75 \\
 & EntropyBal & 3.53$\pm$1.47 & 3.48$\pm$1.52 & 3.56$\pm$1.74 \\
 & AnchorOnly & -- & 4.85$\pm$1.04 & 4.61$\pm$0.87 \\
 & AnchorPlugin & 4.36$\pm$0.98 & 3.95$\pm$1.26 & 4.08$\pm$1.27 \\
 & Proposed & -- & \textbf{0.78$\pm$0.16} & 0.45$\pm$0.12 \\
 & Proposed-CF & -- & 1.16$\pm$0.34 & \textbf{0.39$\pm$0.13} \\
 & Proposed-B & \textbf{3.24$\pm$1.46} & 2.81$\pm$1.43 & 3.26$\pm$1.77 \\
\midrule
10 & TargetOnly & -- & 4.96$\pm$0.81 & 4.66$\pm$0.93 \\
 & ProxyOnly & 5.64$\pm$0.90 & 5.53$\pm$0.94 & 5.41$\pm$1.37 \\
 & IPW-Transport & 3.14$\pm$1.29 & 3.39$\pm$1.39 & 3.46$\pm$1.80 \\
 & OM-Transport & 3.09$\pm$1.30 & 3.28$\pm$1.42 & 3.33$\pm$1.87 \\
 & EntropyBal & 3.31$\pm$1.24 & 3.47$\pm$1.38 & 3.45$\pm$1.81 \\
 & AnchorOnly & -- & 5.04$\pm$0.85 & 4.71$\pm$0.94 \\
 & AnchorPlugin & 4.32$\pm$0.96 & 4.26$\pm$1.11 & 4.21$\pm$1.56 \\
 & Proposed & -- & \textbf{0.73$\pm$0.18} & 0.45$\pm$0.14 \\
 & Proposed-CF & -- & 1.19$\pm$0.39 & \textbf{0.39$\pm$0.15} \\
 & Proposed-B & \textbf{3.09$\pm$1.30} & 2.78$\pm$1.31 & 3.18$\pm$1.85 \\
\midrule
20 & TargetOnly & -- & 4.78$\pm$0.77 & 4.47$\pm$0.83 \\
 & ProxyOnly & 5.71$\pm$0.96 & 5.29$\pm$0.97 & 5.12$\pm$1.00 \\
 & IPW-Transport & 2.86$\pm$1.19 & 2.88$\pm$1.38 & 3.02$\pm$1.56 \\
 & OM-Transport & \textbf{2.84$\pm$1.19} & 2.83$\pm$1.39 & 2.91$\pm$1.58 \\
 & EntropyBal & 2.99$\pm$1.16 & 2.96$\pm$1.36 & 3.03$\pm$1.55 \\
 & AnchorOnly & -- & 4.84$\pm$0.81 & 4.52$\pm$0.83 \\
 & AnchorPlugin & 4.33$\pm$0.97 & 3.92$\pm$1.08 & 3.85$\pm$1.16 \\
 & Proposed & -- & \textbf{0.69$\pm$0.16} & 0.44$\pm$0.10 \\
 & Proposed-CF & -- & 1.08$\pm$0.31 & \textbf{0.39$\pm$0.11} \\
 & Proposed-B & 2.84$\pm$1.19 & 2.39$\pm$1.21 & 2.63$\pm$1.53 \\
\midrule
50 & TargetOnly & -- & 5.05$\pm$1.18 & 4.56$\pm$0.77 \\
 & ProxyOnly & 5.55$\pm$0.85 & 5.62$\pm$1.50 & 5.25$\pm$1.02 \\
 & IPW-Transport & 2.74$\pm$1.17 & 3.23$\pm$1.97 & 2.89$\pm$1.48 \\
 & OM-Transport & 2.74$\pm$1.17 & 3.22$\pm$1.97 & 2.85$\pm$1.50 \\
 & EntropyBal & 2.84$\pm$1.16 & 3.28$\pm$1.94 & 2.92$\pm$1.48 \\
 & AnchorOnly & -- & 5.10$\pm$1.24 & 4.62$\pm$0.78 \\
 & AnchorPlugin & 4.25$\pm$0.90 & 4.27$\pm$1.62 & 3.98$\pm$1.20 \\
 & Proposed & -- & \textbf{0.66$\pm$0.15} & 0.45$\pm$0.16 \\
 & Proposed-CF & -- & 1.26$\pm$0.57 & \textbf{0.41$\pm$0.16} \\
 & Proposed-B & \textbf{2.74$\pm$1.17} & 2.24$\pm$1.61 & 2.38$\pm$1.51 \\
\bottomrule
\end{tabular}
\end{table}

\noindent\textbf{Key findings from the source site scaling sweep:}
\begin{itemize}
\item \textbf{More sources improve all methods}, but the magnitude of improvement varies. Transport baselines (IPW-Transport, OM-Transport) improve from PEHE $\approx 3.5$ at $C=2$ to $\approx 2.7$ at $C=50$.
\item \textbf{\texttt{Proposed} shows particular benefit from source diversity}: with sufficient target budget $(m_0, m_1) = (150, 100)$, PEHE improves from $0.84 \pm 0.16$ at $C=2$ to $0.66 \pm 0.15$ at $C=50$, a 21\% reduction.
\item \textbf{\texttt{Proposed-B} (disconnected target) benefits substantially}: PEHE drops from $3.51 \pm 1.47$ at $C=2$ to $2.74 \pm 1.17$ at $C=50$ in the $(m_0=50, m_1=0)$ regime, matching OM-Transport.
\item \textbf{Diminishing returns}: The largest gains occur when moving from $C=2$ to $C=10$; further increases to $C=50$ provide modest additional improvement, consistent with the automatic source selection filtering out less informative sites.
\end{itemize}

These results demonstrate that the proposed method effectively leverages additional source sites through its automatic detection mechanism, with performance scaling favorably as source diversity increases.

\medskip
\noindent\textbf{Summary and interpretation.}
The source site scaling results directly address the reviewer's concern about site imbalance. Three key observations emerge:

\begin{enumerate}[leftmargin=*]
\item \textbf{Robustness to heterogeneous source pools.} The automatic source detection mechanism in \texttt{Proposed} and \texttt{Proposed-B} guards against negative transfer: as $C$ increases, some sources may be less compatible with the target, but the hypothesis-testing-based screening excludes them. This is evidenced by the monotonic improvement (or at worst, stability) of our methods as $C$ grows, whereas a naive pooling approach would risk degradation from incompatible sources.

\item \textbf{Effective sample size matters more than raw count.} The improvement from $C=2$ to $C=10$ is larger than from $C=10$ to $C=50$, suggesting that once sufficient source diversity is achieved, additional sites contribute diminishing marginal information. This aligns with the theoretical insight that the effective training size $N^\star = \sum_{c \in \widehat{\mathcal{C}}} n_c$ (over \emph{selected} sources) governs finite-sample performance, not the total number of available sites.

\item \textbf{Disconnected targets benefit most from source scaling.} \texttt{Proposed-B} shows the largest relative improvement (22\% PEHE reduction from $C=2$ to $C=50$), consistent with the intuition that when no target treated data is available, the quality and diversity of transportable source CATE information becomes the primary bottleneck. More sources increase the probability of finding compatible ones under Assumption~A6.
\end{enumerate}

\noindent
In summary, the source site scaling experiments confirm that our method gracefully handles varying numbers of heterogeneous source sites, with automatic detection preventing negative transfer while extracting maximal benefit from compatible sources.

\medskip
\noindent\textbf{(ii) Error propagation: why it is second-order in Option~A, and where it remains structural in Option~B.}

\paragraph{Option A (\texttt{Proposed-CF}): orthogonality prevents additive stagewise error accumulation.}
We agree that multi-stage procedures can accumulate error if the final objective is sensitive to nuisance errors.
Our connected-target estimator uses a Neyman-orthogonal DR score with cross-fitting, which makes the final CATE regression \emph{first-order insensitive} to nuisance estimation error. Concretely, letting $\eta$ collect the nuisance functions (here, the arm-specific outcome regressions), the expansion takes the form
\begin{align*}
\widehat\tau_{\mathrm{DR}}(x)-\tau_0(x)
&=
\frac{1}{n_0}\sum_{i:S_i=0}\phi(O_i;x)
\;+\;
R_{n_0}(x),\\
R_{n_0}(x)
&=
O_p\!\Big(
\|\widehat\mu_{0,0}-\mu_{0,0}\|_{L_2(P_0)}
\cdot
\|\widehat\mu_{1,0}-\mu_{1,0}\|_{L_2(P_0)}
\Big)
\;+\;o_p(n_0^{-1/2}),
\end{align*}
so nuisance errors enter through a \emph{product} rather than additively. This is the formal sense in which orthogonalization mitigates error propagation: Stage~1 estimation errors do not linearly propagate into Stage~3, but are attenuated to second order.

\paragraph{Option B (\texttt{Proposed-B}): the dominant issue is structural transport error rather than nuisance propagation.}
In the disconnected regime, the target has placebo only and $\tau_0$ is not identified from target data alone. Our revised procedure does not attempt to estimate a target treated regression via a chain of corrections; instead it screens sources using placebo compatibility and then transports a CATE predictor learned on screened sources. Accordingly, the relevant decomposition is
\begin{align*}
\|\widehat\tau_B-\tau_0\|_{L_2(P_0)}
&\le
\underbrace{\|\widehat\tau_B-\tau^\star\|_{L_2(P_0)}}_{\text{estimation on screened sources}}
\;+\;
\underbrace{\|\tau^\star-\tau_0\|_{L_2(P_0)}}_{\text{structural transport term }(\epsilon_\tau)}
\;+\;
\underbrace{\Delta_{\mathrm{sel}}}_{\text{screening error }(\eta)}.
\end{align*}
Thus, in Option~B the core limitation is not compounding estimation noise across stages, but the irreducible transport approximation term $\epsilon_\tau$ in Assumption~A6 (plus screening failure probability $\eta$). We clarify this distinction in the revision and quantify both terms empirically via the imbalance sweep (which affects $N^\star$ and source diversity) and via controlled ``transport-validity'' sweeps (which affect $\epsilon_\tau$ and $\eta$).

\medskip
\noindent\textbf{Uncertainty quantification and relation to hierarchical/multilevel models.}
We agree that hierarchical/multilevel models provide an alternative that can be attractive under explicit cross-site exchangeability or random-effects assumptions.
However, they typically target a different estimand and require stronger modeling assumptions about between-site variation.
In our setting, we emphasize two complementary uncertainty summaries:
\begin{itemize}
\item In Option~A, we use standard DR-learner inference tools (influence-function based variance for fixed feature maps / linear second stage).
\item In Option~B, uncertainty is primarily governed by (i) estimation variability on screened sources and (ii) sensitivity to the screening step; we therefore report variability across resamples/splits of the source set (and, in experiments, a split-based stability check) as a practical diagnostic.
\end{itemize}

\medskip
\noindent\textbf{Revision summary.}
In the revision we (a) add explicit site-imbalance experiments, (b) make the role of imbalance transparent through the effective training sizes ($n_0$ in Option~A, $N^\star$ in Option~B, and $m_0$ for placebo anchoring when used), and (c) clarify analytically that, in the connected-target regime, orthogonality makes stagewise error propagation second-order (product-form), while in the disconnected-target regime the dominant limitation is structural transport error captured explicitly by Assumption~A6.
}

\clearpage